% Source: https://artofproblemsolving.com/community/c3h3191311_2015_bmmt_individual_round__berkley_mini_math_tournament
% Stable identity: aops-topic-3191311-q18
Take $A=(1,0)$, $C=(-\dfrac{1}{2},\dfrac{\sqrt3}{2})$, $E=(-\dfrac{1}{2},-\dfrac{\sqrt3}{2})$. The other two lines are diameters through $O=(0,0)$ with slopes $\pm\sqrt3$. Intersecting them with $AC$ and $AE$ gives $P=(\dfrac{1}{4},\dfrac{\sqrt3}{4})$ and $Q=(\dfrac{1}{4},-\dfrac{\sqrt3}{4})$. The perpendicular diagonals of the kite have lengths $AO=1$ and $PQ=\dfrac{\sqrt3}{2}$, so the area is $\boxed{\dfrac{\sqrt3}{4}}$.
\begin{asy}
usepackage("amsmath");
size(7cm);
pair[] H;for(int k=0;k<6;++k)H.push(dir(60*k));pair O=(0,0),P=(.25,sqrt(3)/4),Q=(.25,-sqrt(3)/4);draw(circle(O,1));draw(H[4]--H[0]--H[2]);draw(H[5]--H[2]);draw(H[1]--H[4]);filldraw(H[0]--P--O--Q--cycle,lightgray);string[] names={"A","B","C","D","E","F"};for(int k=0;k<6;++k)label("$"+names[k]+"$",H[k],dir(60*k));label("$P$",P,N);label("$Q$",Q,S);label("$O$",O,W);
\end{asy}
